% A155_Meson_Sektor_En_ch.tex — English parallel version of A155
% Meson masses and baryon candidate. Johann Pascher, 23 July 2026

\section{Purpose and Context}

\textbf{[STIPULATION]} This document treats meson masses within the
FFGFT framework using a self-contained additive ansatz. It establishes
the ansatz's domain of validity, evaluates it on the pion, documents
the connection to the Z$_3$ sector structure (A270), and states the
structural boundary at which the additive ansatz ends.

A150 treats the quark sector as descriptive classification, not as a
predictive quantity. The same caution applies here: quark masses are
scheme-dependent (A150, \S{}Quarks, first reason); all numerical values
use PDG figures ($\overline{\text{MS}}$, 2\,GeV) and carry their
uncertainty ($\sim$1--2\,\%).

\section{External Inputs}

\textbf{[STIPULATION]} Besides the A-series quantities $\xi$ and
$K_\text{frak} = 1-100\xi$ (derivation: A040), this document uses three
external standard quantities of QCD; they are declared here, not
derived:

\begin{center}
\renewcommand{\arraystretch}{1.3}
{\small\begin{tabular}{llp{6cm}}
\toprule
Quantity & Value & Role \\
\midrule
$\Lambda_\text{QCD}$ & $0{.}217$\,GeV & QCD scale (PDG range
$0{.}21$--$0{.}22$\,GeV); sole dimensionful input of the binding term \\
$m_u, m_d, m_s$ & $2{.}16$; $4{.}67$; $93{.}4$\,MeV &
current quark masses ($\overline{\text{MS}}$, 2\,GeV, PDG) \\
Meson masses & PDG & comparison values \\
\bottomrule
\end{tabular}}
\renewcommand{\arraystretch}{1.0}
\end{center}

\section{The Additive Ansatz}

\textbf{[STIPULATION]} For a meson of two constituents we set:
\[
  m_M = m_{q_1} + m_{q_2}
  + \Lambda_\text{QCD}\cdot K_\text{frak}^{\,n_\text{eff}}.
\]
Rationale for the form: constituent masses enter additively; the
binding part is taken as the QCD scale times a fractal attenuation
$K_\text{frak}^{\,n_\text{eff}}$ --- the same $K_\text{frak}$ power
structure that carries scale transitions throughout the corpus (A040;
level overview P40, Doc.~190). The exponent $n_\text{eff}$ is initially
free and is evaluated against experiment; its topological
interpretation is supplied by A270.

\section{Domain of Validity}

\textbf{[K]} Since $K_\text{frak} < 1$, for every $n_\text{eff} > 0$:
\[
  \Lambda_\text{QCD}\cdot K_\text{frak}^{\,n_\text{eff}} < \Lambda_\text{QCD}.
\]
\textbf{The additive ansatz can structurally carry only binding
contributions below $\Lambda_\text{QCD}$.} This is a property of the
functional form, not an approximation issue.

\medskip
\textbf{[K]} Numerical check:

\begin{center}
\renewcommand{\arraystretch}{1.3}
{\small\begin{tabular}{lcccc l}
\toprule
Meson & $m_M$ (GeV) & $\sum m_q$ & Binding & $n_\text{eff}$ & Status \\
\midrule
$\pi^0$ ($u\bar u/d\bar d$) & 0{.}1350 & 0{.}0043 & 0{.}1307 & 37{.}79 & within domain \\
$\pi^\pm$ ($u\bar d$)       & 0{.}1396 & 0{.}0068 & 0{.}1327 & 36{.}62 & within domain \\
$K^\pm$ ($u\bar s$)         & 0{.}4937 & 0{.}0956 & 0{.}3981 & --- & binding $>\Lambda$ \\
$\eta$                      & 0{.}5479 & 0{.}1868 & 0{.}3611 & --- & binding $>\Lambda$ \\
$\rho$ ($u\bar u$)          & 0{.}7753 & 0{.}0043 & 0{.}7709 & --- & binding $>\Lambda$ \\
$\omega$ ($u\bar u$)        & 0{.}7827 & 0{.}0043 & 0{.}7783 & --- & binding $>\Lambda$ \\
\bottomrule
\end{tabular}}
\renewcommand{\arraystretch}{1.0}
\end{center}

Only the pions lie within the domain. The failure for $K$, $\eta$,
$\rho$, $\omega$ is not a numerical defect but the declared boundary of
the functional form.

\section{Sector Connection (A270)}

\textbf{[K]} From the Z$_3$ sector structure, A270 predicts for mesons
(1~twisted sector) the exponent $36+1 = 37$; the exact reference value
of the bulk relation is $n = 35{.}99$
($K_\text{frak}^{-36} \approx 16/\pi^2$, A270 \S{}Key Relation).

Evaluating the pion data yields:
\[
  n_\text{eff}(\pi^0) = 37{.}79, \qquad n_\text{eff}(\pi^\pm) = 36{.}62.
\]
Both values lie within the window of the sector prediction $36+1$.
\textbf{The meson step of the A270 sector table is thereby empirically
anchored at the pion} --- the exponent was determined here from the
mass data independently of the sector argument.

\section{GMOR Check and the Limit of the Additive Ansatz}

\textbf{[K]} For pseudo-Goldstone bosons the Gell-Mann--Oakes--Renner
relation of chiral perturbation theory holds (external standard
physics):
\[
  m_P^2 = (m_{q_1} + m_{q_2})\cdot B,
  \qquad B = \frac{-2\langle\bar qq\rangle}{f_\pi^2}.
\]
Evaluation with the external inputs:
\[
  B_\pi = \frac{m_{\pi^\pm}^2}{m_u+m_d} = 2{.}852\ \text{GeV},
  \qquad
  B_K = \frac{m_K^2}{m_u+m_s} = 2{.}550\ \text{GeV},
  \qquad
  \frac{B_\pi}{B_K} = 1{.}12.
\]
Pion and kaon share the same chiral structure ($B$ constant up to the
known SU(3) breaking). \textbf{The additive ansatz thus fails for the
kaon not because of different physics but because of the wrong
functional form}: linear instead of quadratic (GMOR).

\medskip
\textbf{[S]} \textbf{Candidate: $B$ carries the sector factor.}
Numerically one notes:
\[
  \frac{B_\pi}{\Lambda_\text{QCD}\cdot K_\text{frak}^{-37}}
  = 7{.}999 \approx 8,
  \qquad\text{i.e.}\qquad
  B \approx 8\,\Lambda_\text{QCD}\cdot K_\text{frak}^{-37}.
\]
The~8 admits two corpus readings: $\dim\text{SU(3)} = 8$ (number of
gluons --- the chiral condensate is gluonically carried) or the D4
ratio Weyl order/kissing number $= 192/24 = 8$ (A270).
\textbf{P35 applies}: the central deviation ($0{.}02\,\%$) lies far
below the scheme uncertainty of the quark masses ($\sim$2\,\%); the hit
is within uncertainty, but the precision is not robust. Booked as
candidate, not as result.

Were the candidate viable, a GMOR form with sector factor would follow,
\[
  m_P^2 = (m_{q_1}+m_{q_2})\cdot 8\,\Lambda_\text{QCD}
  \cdot K_\text{frak}^{-37},
\]
covering pion \emph{and} kaon (kaon then to $\sim$6\,\% via SU(3)
breaking in $B$). A [K]-level test of this form requires scheme-fixed
quark masses --- open.

\medskip
\textbf{[STIPULATION]} \textbf{Vector mesons} ($\rho$, $\omega$) are
not Goldstone bosons; neither the additive ansatz nor GMOR applies to
them. Their masses sit at the chiral symmetry-breaking scale itself
($\sim 4\pi f_\pi$). An FFGFT treatment of the vector sector does not
exist and is not claimed here.

\section{Baryon Candidate: Fully Geometric Proton Formula}

\textbf{[S]} The additive meson ansatz and the sector structure (A270)
raise the question whether the proton mass, too, admits a geometric
decomposition. Candidate:
\[
  m_p = \underbrace{\frac{\pi^2}{2}}_{\text{Vol}(B^4)}
  \cdot \underbrace{\frac{\pi}{6}}_{\eta_\text{sc}}
  \cdot \underbrace{N_c\,(1+\alpha_s)}_{\text{color/coupling}}
  \cdot \underbrace{e^{-3\xi/4}}_{\xi\text{-damping}}
  \cdot \underbrace{\frac{\Lambda_\text{QCD}}{2}}_{\text{scale}}
  = 0{.}9402\ \text{GeV}
\]
against PDG $0{.}93827$ ($\Delta = +0{.}21\,\%$). Every factor is
interpreted:

\begin{center}
\renewcommand{\arraystretch}{1.3}
{\small\begin{tabular}{llp{6.2cm}}
\toprule
Factor & Value & Interpretation \\
\midrule
$\pi^2/2$ & $4{.}9348$ & volume of the unit 4-ball (T$^4$-native) \\
$\pi/6$ & $0{.}5236$ & packing density of the simple cubic lattice,
$\text{Vol}(B^3)/2^3$ --- 3D analogue of the D4 density $\pi^2/16$
(4D) that carries the bulk exponent (A270) \\
$N_c(1+\alpha_s)$ & $3{.}354$ & color number times coupling correction
(external QCD quantities, \S{}External Inputs) \\
$e^{-3\xi/4}$ & $0{.}9999$ & $\xi$-damping ($N_c$-fold) \\
$\Lambda_\text{QCD}/2$ & $0{.}1085$\,GeV & halved QCD scale; sole
dimensionful carrier \\
\bottomrule
\end{tabular}}
\renewcommand{\arraystretch}{1.0}
\end{center}

\textbf{Absorbability of the residual:} $+0{.}21\,\%$ closes exactly
with $\alpha_s = 0{.}1156$ (instead of $0{.}118$; $-2\,\%$) or
$\Lambda_\text{QCD} = 0{.}2165$ (instead of $0{.}217$; $-0{.}2\,\%$)
--- both deep within the uncertainty of the external inputs.

\medskip
\textbf{[S]} \textbf{Structural link to the Higgs formula.} The two
geometry factors combine to
\[
  \frac{\pi^2}{2}\cdot\frac{\pi}{6} = \frac{\pi^3}{12},
  \qquad\text{and}\qquad
  \frac{16\pi^3}{\pi^3/12} = 192
\]
--- the D4 Weyl order. Proton geometry and the Higgs formula $16\pi^3$
(Doc.~190, K1) thus stand in the exact ratio of the Weyl number 192.
Noted under P35, not booked.

\medskip
\textbf{P35 caveat (competing candidate):} The product of free factors
fixed by the calculation is $X = 0{.}52248$; $\pi/6 = 0{.}52360$ hits
at $+0{.}21\,\%$. Numerically closer would be
$K_\text{frak}^{48{.}4}$ ($-0{.}05\,\%$), but without integer exponent,
without interpretation, and without an absorbability argument
($K_\text{frak}^{48}$ misses at $+0{.}49\,\%$). Two candidates, no
forward derivation: [S], not [K]. The historical occasion of this
analysis (correction of a legacy-corpus calculation) is registered in
Doc.~190, R62.

\section{Layer-Status Summary}

\begin{center}
\renewcommand{\arraystretch}{1.3}
{\small\begin{tabular}{ll}
\toprule
Statement & Status \\
\midrule
Additive ansatz $m_M = \sum m_q + \Lambda\,K_\text{frak}^{n}$ & [STIPULATION] (ansatz) \\
External inputs $\Lambda_\text{QCD}$, quark masses & [STIPULATION] (declared) \\
Domain: only binding $<\Lambda_\text{QCD}$ & [K] (form property) \\
Pion within domain, $n_\text{eff}(\pi^0)=37{.}79$ & [K] \\
Sector connection $n_\text{eff}\approx 37 = 36+1$ (A270) & [K] \\
GMOR consistency $B_\pi \approx B_K$ & [K] (standard ChPT) \\
$B \approx 8\,\Lambda\cdot K_\text{frak}^{-37}$ & [S] (candidate, P35) \\
Vector mesons outside & [STIPULATION] \\
Proton formula $\frac{\pi^3}{12}N_c(1+\alpha_s)e^{-3\xi/4}\frac{\Lambda}{2}$ & [S] (candidate, P35) \\
$16\pi^3/(\pi^3/12)=192$ (Higgs link) & [S] (noted, P35) \\
\bottomrule
\end{tabular}}
\renewcommand{\arraystretch}{1.0}
\end{center}

\section{Open Points}

\begin{enumerate}
\item Scheme-fixed quark masses for a [K]-level test of the
  $8\Lambda K_\text{frak}^{-37}$ candidate.
\item Origin of the~8: gluon number or D4 ratio --- both readings
  unconnected; unbooked under P35 until a forward derivation exists.
\item $\pi^\pm$/$\pi^0$ difference ($36{.}62$ vs.\ $37{.}79$): does the
  electromagnetic self-energy of the charged pion carry a sector
  contribution of its own?
\item Vector meson sector: fully open.
\item Forward derivation of the baryon candidate: why $\pi/6$ (3D
  packing) alongside $\text{Vol}(B^4)$ (4D volume)? Decision between
  the candidates $\pi/6$ and a $K_\text{frak}$ power pending.
\item For the correction of the historical legacy-corpus baryon
  calculation see Doc.~190, register R62.
\end{enumerate}

\section{Notation}

Symbols used throughout the A-series are explained in A010.
Specific to this document:

\begin{center}
\renewcommand{\arraystretch}{1.3}
{\small\begin{tabular}{lp{9cm}}
\toprule
Symbol & Meaning \\
\midrule
$m_M$, $m_P$ & meson mass in general resp.\ pseudoscalar \\
$m_{q_i}$ & current quark masses ($\overline{\text{MS}}$, 2\,GeV, PDG; \S{}External Inputs) \\
$\Lambda_\text{QCD}=0{.}217$\,GeV & QCD scale (\S{}External Inputs) \\
$n_\text{eff}$ & exponent of the fractal binding correction $K_\text{frak}^{\,n_\text{eff}}$ \\
$B$ & GMOR parameter, $B=-2\langle\bar qq\rangle/f_\pi^2$; from $m_P^2=(m_{q_1}+m_{q_2})\,B$ \\
$\langle\bar qq\rangle$, $f_\pi$ & chiral condensate; pion decay constant \\
$\eta_\text{sc}=\pi/6$ & packing density of the simple cubic lattice \\
$\text{Vol}(B^4)=\pi^2/2$ & volume of the unit 4-ball \\
$N_c=3$, $\alpha_s=0{.}118$ & color number; strong coupling (external QCD quantities) \\
\bottomrule
\end{tabular}}
\renewcommand{\arraystretch}{1.0}
\end{center}

\medskip\noindent\textbf{Verification script:} \texttt{python/A\_Serie\_Skripte/a155\_meson.py}
--- all numerical statements of this document are checked there with assertions.

\section{References}

\begin{center}
{\small\begin{tabular}{ll}
\toprule
Document & Reference \\
\midrule
A010 & A-series notation \\
A040 & derivation of $K_\text{frak} = 1-100\xi$ \\
A110 & circulant (Z$_3$ algebra of the untwisted sector) \\
A150 & quark-sector caution: scheme dependence, no predictive quantity \\
A270 & Z$_3$ sector structure, meson exponent $36+1=37$ \\
Doc.~190 & P35, P40; register R62 (baryon correction, legacy corpus) \\
\bottomrule
\end{tabular}}
\end{center}
